\subsection{Operadores diferenciales.}
\footnote{Para un formalismo más general, véase la Exposición VII de [42]
(debida a P. Gabriel).}
\label{IV.16.8}

\begin{definition}[16.8.1]
\label{IV.16.8.1}
Sean $f=(\psi,\theta):X\to S$ un morfismo de preesquemas,
$\mathscr{F}$, $\mathscr{G}$ dos $\mathcal{O}_X$-módulos, y $n$ un entero
$\geq0$. Se dice que un homomorfismo de haces de grupos aditivos
$D:\mathscr{F}\to\mathscr{G}$ es un \emph{operador diferencial de orden
$\leq n$} (relativamente a $S$) si existe un homomorfismo de
$\mathcal{O}_X$-módulos
$u:\mathscr{P}_{X/S}^n(\mathscr{F})\to\mathscr{G}$ (donde
$\mathscr{P}_{X/S}^n(\mathscr{F})$ está provisto de su estructura de
$\mathcal{O}_X$-módulo a la izquierda (16.7.4)) tal que
$D=u\circ d_{X/S,\mathscr{F}}^n$.
\end{definition}

Es claro, en virtud de la existencia de los homomorfismos canónicos
\[
  \mathscr{P}_{X/S}^m(\mathscr{F})
  \longrightarrow\mathscr{P}_{X/S}^n(\mathscr{F})
\]
\oldpage[IV-4]{40}
para $n\leq m$ (16.7.7), que un operador diferencial de orden $\leq n$ es
también un operador diferencial de orden $\leq m$ para todo $m\geq n$. Si
$D:\mathscr{F}\to\mathscr{G}$ es un operador diferencial de orden
$\leq n$, entonces, para todo abierto $U$ de $X$,
$D|U:\mathscr{F}|U\to\mathscr{G}|U$ es también un operador diferencial
de orden $\leq n$.

Se dice que un homomorfismo $D:\mathscr{F}\to\mathscr{G}$ de los haces de
grupos aditivos subyacentes a $\mathscr{F}$ y $\mathscr{G}$ es un
\emph{operador diferencial} (relativamente a $S$) si, para todo $x\in X$,
existen una vecindad abierta $U$ de $x$ y un entero $n\geq0$ tales que
$D|U:\mathscr{F}|U\to\mathscr{G}|U$ sea un operador diferencial de orden
$\leq n$. El \emph{orden} de un operador diferencial
$D:\mathscr{F}\to\mathscr{G}$ es el ínfimo de los enteros $n$ tales que
$D$ sea un operador diferencial de orden $\leq n$ (y, por tanto, $+\infty$
si no hay tales enteros); este orden es siempre finito si $X$ es
\emph{cuasicompacto}. Los operadores diferenciales de orden $0$ no son sino
los homomorfismos de $\mathcal{O}_X$-módulos
$\mathscr{F}\to\mathscr{G}$; se conviene en que todo operador diferencial
de orden $<0$ es \emph{nulo}. Para $n\geq0$, un operador diferencial no es,
en general, un homomorfismo de $\mathcal{O}_X$-módulos, pero es siempre
un homomorfismo de $\psi^*(\mathcal{O}_S)$-módulos.

Cuando $\mathscr{F}=\mathcal{O}_X$, un operador diferencial de orden
$\leq1$ de $\mathcal{O}_X$ en $\mathscr{G}$ se escribe de manera única
en la forma $v+D$, donde $v:\mathcal{O}_X\to\mathscr{G}$ es un
$\mathcal{O}_X$-homomorfismo, y $D$ una $S$-\emph{derivación} (16.5.1) de
$\mathcal{O}_X$ en $\mathscr{G}$: esto resulta de la estructura de
$P_{B/A}^1$ (0, 20.4.8).

\begin{env}[16.8.2]
\label{IV.16.8.2}
Para describir de manera más explícita un operador diferencial de orden
$\leq n$, $D:\mathscr{F}\to\mathscr{G}$, basta, para todo abierto afín
$U$ de $X$ cuya imagen en $S$ esté contenida en un abierto afín $V$,
caracterizar el homomorfismo
$D=D_U:\Gamma(U,\mathscr{F})\to\Gamma(U,\mathscr{G})$. Si se pone
$\Gamma(V,\mathcal{O}_S)=A$, $\Gamma(U,\mathcal{O}_X)=B$, de modo que $B$
es una $A$-álgebra, se tiene
\[
  \Gamma(U,\mathscr{P}_{X/S}^n)
  =(B\otimes_A B)/\mathfrak{I}^{n+1},
\]
donde se ha puesto, para abreviar, $\mathfrak{I}=\mathfrak{I}_{B/A}$.
Pongamos además $M=\Gamma(U,\mathscr{F})$, $N=\Gamma(U,\mathscr{G})$;
entonces la definición de $D$ significa que, para cada par $(U,V)$ que
satisface las condiciones precedentes, el $A$-homomorfismo $D:M\to N$ se
factoriza como
\[
  M\longrightarrow
  ((B\otimes_A B)/\mathfrak{I}^{n+1})\otimes_B M
  \xrightarrow{v}N
\]
donde la primera flecha es el homomorfismo canónico $t\mapsto1\otimes t$, y
$v$ es un $B$-homomorfismo; la estructura de $B$-módulo de
$((B\otimes_A B)/\mathfrak{I}^{n+1})\otimes_B M$ proviene del primer
factor $B$ (mientras que, recordemos, al formar el producto tensorial
sobre $B$, la estructura de $B$-módulo de
$(B\otimes_A B)/\mathfrak{I}^{n+1}$ proviene del segundo factor $B$).
Observemos ahora que el $B$-módulo
$((B\otimes_A B)/\mathfrak{I}^{n+1})\otimes_B M$ es isomorfo a
$(B\otimes_A M)/\mathfrak{I}^{n+1}(B\otimes_A M)$, donde $B\otimes_A M$ se
considera como $(B\otimes_A B)$-módulo y su estructura de $B$-módulo
proviene del homomorfismo $b\mapsto b\otimes1$ de $B$ en $B\otimes_A B$.
Sea entonces $D'$ el $B$-homomorfismo de $B\otimes_A M$ en $N$ tal que
$D'(b\otimes t)=bD(t)$; la condición de factorización de $D$ se expresa
también diciendo que $D'$ debe \emph{anularse sobre el $B$-submódulo}
$\mathfrak{I}^{n+1}(B\otimes_A M)$.
\end{env}

\begin{env}[16.8.3]
\label{IV.16.8.3}
Es claro que el conjunto de los operadores diferenciales de orden $\leq n$ de
$\mathscr{F}$ en $\mathscr{G}$ es un grupo aditivo, denotado por
$\operatorname{Diff}_{X/S}^n(\mathscr{F},\mathscr{G})$; cuando
$\mathscr{F}=\mathscr{G}=\mathcal{O}_X$, se escribe también
$\operatorname{Diff}_{X/S}^n$ en lugar de
$\operatorname{Diff}_{X/S}^n(\mathcal{O}_X,\mathcal{O}_X)$.

Se ha visto en (16.8.1) que, para dos abiertos $U\supset V$ de $X$, se tiene
un homomorfismo canónico de restricción
\[
  \operatorname{Diff}_{U/S}^n(\mathscr{F}|U,\mathscr{G}|U)
  \longrightarrow
  \operatorname{Diff}_{V/S}^n(\mathscr{F}|V,\mathscr{G}|V)
\]
\oldpage[IV-4]{41}
por lo que $U\mapsto\operatorname{Diff}_{U/S}^n(\mathscr{F}|U,\mathscr{G}|U)$
es un prehaz de grupos aditivos; de hecho, es incluso un \emph{haz}, pues,
cuando $U$ recorre los abiertos de $X$, los homomorfismos
$u\mapsto u\circ d_{U/S,\mathscr{F}|U}^n$ son \emph{isomorfismos} de grupos
aditivos
\[
  \mathscr{H}\!om_{\mathcal{O}_U}
  (\mathscr{P}_{U/S}^n(\mathscr{F}|U),\mathscr{G}|U)
  \xrightarrow{\sim}
  \operatorname{Diff}_{U/S}^n(\mathscr{F}|U,\mathscr{G}|U),
  \tag{16.8.3.1}
  \label{IV.16.8.3.1}
\]
en virtud de que la imagen de $\mathscr{F}$ por
$d_{X/S,\mathscr{F}}^n$ engendra $\mathscr{P}_{X/S}^n(\mathscr{F})$
(16.7.6). Se denota este haz por
$\mathscr{D}\!iff_{X/S}^n(\mathscr{F},\mathscr{G})$, y se tiene, por tanto:
\end{env}

\begin{proposition}[16.8.4]
\label{IV.16.8.4}
Los isomorfismos (16.8.3.1) definen un isomorfismo de haces de grupos
aditivos
\[
  \mathscr{H}\!om_{\mathcal{O}_X}
  (\mathscr{P}_{X/S}^n(\mathscr{F}),\mathscr{G})
  \xrightarrow{\sim}
  \mathscr{D}\!iff_{X/S}^n(\mathscr{F},\mathscr{G}).
  \tag{16.8.4.1}
  \label{IV.16.8.4.1}
\]
\end{proposition}

Cuando $\mathscr{F}=\mathscr{G}=\mathcal{O}_X$, se escribe también
$\mathscr{D}\!iff_{X/S}^n$ en lugar de
$\mathscr{D}\!iff_{X/S}^n(\mathcal{O}_X,\mathcal{O}_X)$; de (16.8.4) resulta
que $\mathscr{D}\!iff_{X/S}^n$ se identifica canónicamente con el
\emph{dual} del $\mathcal{O}_X$-módulo $\mathscr{P}_{X/S}^n$; por ello se
escribe $\langle t,D\rangle$ en lugar de $u(t)$ si $t$ es una sección de
$\mathscr{P}_{X/S}^n$ sobre un abierto y si $u$ es el homomorfismo de
$\mathscr{P}_{X/S}^n$ en $\mathcal{O}_X$ correspondiente a $D$.

\begin{env}[16.8.5]
\label{IV.16.8.5}
Puesto que $\mathscr{P}_{X/S}^n(\mathscr{F})$ está provisto de una estructura
de $\mathcal{O}_X$-bimódulo (16.7.4), se deduce canónicamente una estructura
de $\mathcal{O}_X$-bimódulo sobre
$\mathscr{H}\!om_{\mathcal{O}_X}
(\mathscr{P}_{X/S}^n(\mathscr{F}),\mathscr{G})$, y, por tanto, también sobre
$\mathscr{D}\!iff_{X/S}^n(\mathscr{F},\mathscr{G})$ en virtud de (16.8.4.1).
Más precisamente, a la estructura de $\mathcal{O}_X$-módulo a la izquierda
de $\mathscr{P}_{X/S}^n(\mathscr{F})$ corresponde, en virtud de la definición
(16.8.1), la estructura de $\mathcal{O}_X$-módulo a la izquierda sobre
$\mathscr{D}\!iff_{X/S}^n(\mathscr{F},\mathscr{G})$ que se explicita como
sigue: para todo abierto $U$ de $X$, toda sección
$a\in\Gamma(U,\mathcal{O}_X)$ y todo operador diferencial
$D:\mathscr{F}|U\to\mathscr{G}|U$, $aD$ es el operador diferencial que a
toda sección $t\in\Gamma(U,\mathscr{F})$ hace corresponder la sección
\[
  (aD)(t)=a(D(t))
  \tag{16.8.5.1}
  \label{IV.16.8.5.1}
\]
de $\Gamma(U,\mathscr{G})$. Análogamente, a la estructura de
$\mathcal{O}_X$-módulo a la derecha de $\mathscr{P}_{X/S}^n(\mathscr{F})$
corresponde la estructura de $\mathcal{O}_X$-módulo a la derecha sobre
$\mathscr{D}\!iff_{X/S}^n(\mathscr{F},\mathscr{G})$ que se explicita como
sigue: con las mismas notaciones que antes, $Da$ es el operador diferencial
que a $t\in\Gamma(U,\mathscr{F})$ hace corresponder la sección
\[
  (Da)(t)=D(at).
  \tag{16.8.5.2}
  \label{IV.16.8.5.2}
\]
\end{env}

\begin{proposition}[16.8.6]
\label{IV.16.8.6}
Si $f:X\to S$ es un morfismo localmente de presentación finita,
$\mathscr{F}$ un $\mathcal{O}_X$-módulo cuasicoherente de presentación finita
y $\mathscr{G}$ un $\mathcal{O}_X$-módulo cuasicoherente, entonces
$\mathscr{D}\!iff_{X/S}^n(\mathscr{F},\mathscr{G})$ es un
$\mathcal{O}_X$-módulo cuasicoherente para una u otra de las estructuras
definidas en (16.8.5).
\end{proposition}

La proposición resulta de que, bajo las hipótesis formuladas,
$\mathscr{P}_{X/S}^n(\mathscr{F})$ es un $\mathcal{O}_X$-módulo
cuasicoherente de presentación finita (16.7.4), y de (I, 1.3.12).

\begin{env}[16.8.7]
\label{IV.16.8.7}
El conjunto de los operadores diferenciales de $\mathscr{F}$ en
$\mathscr{G}$ (de orden no especificado (16.8.1)) se denota por
$\operatorname{Diff}_{X/S}(\mathscr{F},\mathscr{G})$; se ve también, como en
(16.8.3), que
$U\mapsto\operatorname{Diff}_{U/S}(\mathscr{F}|U,\mathscr{G}|U)$ es un haz de
grupos aditivos, que se designará por
$\mathscr{D}\!iff_{X/S}(\mathscr{F},\mathscr{G})$. Es inmediato que
$\mathscr{D}\!iff_{X/S}(\mathscr{F},\mathscr{G})$ es la unión de la familia
filtrante creciente de sus subhaces
$\mathscr{D}\!iff_{X/S}^n(\mathscr{F},\mathscr{G})$; si $X$ es
cuasicompacto, $\operatorname{Diff}_{X/S}(\mathscr{F},\mathscr{G})$
\oldpage[IV-4]{42}
es asimismo la unión de sus subgrupos
$\operatorname{Diff}_{X/S}^n(\mathscr{F},\mathscr{G})$ (16.8.1). Las
estructuras de $\mathcal{O}_X$-bimódulo sobre los
$\mathscr{D}\!iff_{X/S}^n(\mathscr{F},\mathscr{G})$ definen, por tanto, una
estructura de $\mathcal{O}_X$-bimódulo sobre
$\mathscr{D}\!iff_{X/S}(\mathscr{F},\mathscr{G})$, explicitada también por
(16.8.5.1) y (16.8.5.2).

Obsérvese que, para $n\leq m$, se tiene un diagrama conmutativo
\[
  \xymatrix{
    \mathscr{H}\!om_{\mathcal{O}_X}
      (\mathscr{P}_{X/S}^n(\mathscr{F}),\mathscr{G})
      \ar[r]^-{\sim}\ar[d]
      &\mathscr{D}\!iff_{X/S}^n(\mathscr{F},\mathscr{G})\ar[d]\\
    \mathscr{H}\!om_{\mathcal{O}_X}
      (\mathscr{P}_{X/S}^m(\mathscr{F}),\mathscr{G})
      \ar[r]^-{\sim}
      &\mathscr{D}\!iff_{X/S}^m(\mathscr{F},\mathscr{G})
  }
  \tag{16.8.7.1}
  \label{IV.16.8.7.1}
\]
donde las flechas horizontales son los isomorfismos
(16.8.4.1)\footnote{El texto impreso remite a (16.8.4.2), que no existe; la
remisión correcta es a (16.8.4.1).} y la flecha vertical de la izquierda
proviene del homomorfismo canónico
$\mathscr{P}_{X/S}^m(\mathscr{F})\to
\mathscr{P}_{X/S}^n(\mathscr{F})$ (16.7.7). Para todo abierto $U$ de $X$,
provistamos entonces
\[
  \Gamma(U,\mathscr{P}_{X/S}^{\infty}(\mathscr{F}))
  =\varprojlim\Gamma(U,\mathscr{P}_{X/S}^n(\mathscr{F}))
\]
de la topología de límite proyectivo de las topologías discretas sobre los
$\Gamma(U,\mathscr{P}_{X/S}^n(\mathscr{F}))$, lo que define
$\Gamma(U,\mathscr{P}_{X/S}^{\infty}(\mathscr{F}))$ como un
$\Gamma(U,\mathcal{O}_X)$-bimódulo topológico, de modo que
$\mathscr{P}_{X/S}^{\infty}(\mathscr{F})$ aparece como un haz con valores en
la categoría de los grupos conmutativos topológicos
(0\textsubscript{I}, 3.2.6). Entonces, por (G, II, 1.11), el límite del sistema
inductivo de haces de grupos conmutativos
\[
  (\mathscr{H}\!om_{\mathcal{O}_X}
  (\mathscr{P}_{X/S}^n(\mathscr{F}),\mathscr{G}))
\]
no es sino el haz de gérmenes de homomorfismos \emph{continuos} de
$\mathscr{P}_{X/S}^{\infty}(\mathscr{F})$ en $\mathscr{G}$ (este último
provisto de la topología discreta): en efecto, los homomorfismos continuos de
$\Gamma(U,\mathscr{P}_{X/S}^{\infty}(\mathscr{F}))$ en el grupo discreto
$\Gamma(U,\mathscr{G})$ corresponden biunívocamente a los sistemas inductivos
de homomorfismos de grupos
$\Gamma(U,\mathscr{P}_{X/S}^n(\mathscr{F}))\to\Gamma(U,\mathscr{G})$. Se
puede, por tanto, expresar también (16.8.4) diciendo que se tiene un
isomorfismo canónico
\[
  \mathscr{H}\!om.\mathrm{cont}_{\mathcal{O}_X}
  (\mathscr{P}_{X/S}^{\infty}(\mathscr{F}),\mathscr{G})
  \xrightarrow{\sim}
  \mathscr{D}\!iff_{X/S}(\mathscr{F},\mathscr{G})
\]
donde el primer miembro designa el haz de gérmenes de homomorfismos continuos
de $\mathscr{P}_{X/S}^{\infty}(\mathscr{F})$ en $\mathscr{G}$.
\end{env}

\begin{proposition}[16.8.8]
\label{IV.16.8.8}
Sean $\mathscr{F}$, $\mathscr{G}$ dos $\mathcal{O}_X$-módulos,
$D:\mathscr{F}\to\mathscr{G}$ un homomorfismo de
$\psi^*(\mathcal{O}_S)$-módulos, y $n$ un entero $\geq0$. Las condiciones
siguientes son equivalentes:
\begin{enumerate}
  \item[(a)] $D$ es un operador diferencial de orden $\leq n$.
  \item[(b)] Para toda sección $a$ de $\mathcal{O}_X$ sobre un abierto
  $U$, el homomorfismo $D_a:\mathscr{F}|U\to\mathscr{G}|U$ tal que, para
  toda sección $t$ de $\mathscr{F}$ sobre un abierto $V\subset U$, se tenga
  \[
    D_a(t)=D(at)-aD(t)
    \tag{16.8.8.1}
    \label{IV.16.8.8.1}
  \]
  es un operador diferencial de orden $\leq n-1$.
  \item[(c)] Para todo abierto $U$ de $X$, toda familia
  $(a_i)_{1\leq i\leq n+1}$ de $n+1$ secciones de $\mathcal{O}_X$ sobre
  $U$ y toda sección $t$ de $\mathscr{F}$ sobre $U$, se tiene la identidad
  \[
    \sum_{H\subset I_{n+1}}(-1)^{\operatorname{Card}(H)}
    \left(\prod_{i\in H}a_i\right)
    D\!\left(\left(\prod_{i\notin H}a_i\right)t\right)=0
    \tag{16.8.8.2}
    \label{IV.16.8.8.2}
  \]
  (donde $I_{n+1}$ es el intervalo $1\leq i\leq n+1$ de $\mathbb{N}$).
\end{enumerate}
\oldpage[IV-4]{43}
\end{proposition}

Demostremos primero la equivalencia entre a) y c). Por definición, para
demostrar que $D$ es un operador diferencial de orden $\leq n$, basta
demostrarlo para la restricción
$D|U:\mathscr{F}|U\to\mathscr{G}|U$ a todo abierto afín $U$ de $X$; por
otra parte, la propiedad c) es válida para todo abierto $U$ de $X$ si lo es
para todo abierto afín. Se puede, por tanto, reducir al caso en que
$S=\operatorname{Spec}(A)$ y $X=\operatorname{Spec}(B)$ son afines. En
virtud de (16.8.2) (cuyas notaciones se conservan), la condición a) significa
entonces que el $A$-homomorfismo $D':B\otimes_A M\to N$ tal que
$D'(b\otimes t)=bD(t)$ se anula sobre
$\mathfrak{I}^{n+1}(B\otimes_A M)$, lo que, en virtud de (0, 20.4.4),
equivale a decir que $D'$ se anula para todos los elementos de la forma
\[
  \left(\prod_{i=1}^{n+1}(a_i\otimes1-1\otimes a_i)\right).(1\otimes t)
\]
donde $a_i\in B$ y $t\in M$. Ahora bien, este elemento se escribe, salvo el
signo común $(-1)^{n+1}$, como\footnote{El texto impreso omite en esta
expansión los signos $(-1)^{\operatorname{Card}(H)}$; sin ellos, la imagen
por $D'$ no es el primer miembro de (16.8.8.2).}
\[
  \sum_{H\subset I_{n+1}}(-1)^{\operatorname{Card}(H)}
  \left(\prod_{i\in H}a_i\right)\otimes
  \left(\left(\prod_{i\notin H}a_i\right)t\right),
\]
y el valor de $D'$ sobre este elemento no es sino, salvo ese signo común, el
primer miembro de (16.8.8.2), lo que demuestra la equivalencia entre a) y c).

Demostremos ahora la equivalencia entre b) y c). Razonemos por inducción
sobre $n$, siendo trivial la afirmación para $n=0$. Escribiendo $a_{n+1}$ en
lugar de $a$ en la condición b), se ve, en virtud de la hipótesis de
inducción, que la condición b) significa que, para toda familia
$(a_i)_{1\leq i\leq n}$ de $n$ secciones de $\mathcal{O}_X$ sobre $U$ y
toda sección $t$ de $\mathscr{F}$ sobre $U$, se tiene
\[
  \sum_{H'\subset I_n}(-1)^{\operatorname{Card}(H')}
  \left(\prod_{i\in H'}a_i\right)
  D_{a_{n+1}}\!\left(\left(\prod_{i\notin H'}a_i\right)t\right)=0.
\]
Pero, si en esta relación se sustituye $D_{a_{n+1}}$ por su definición
(16.8.8.1), se comprueba inmediatamente que se obtiene, salvo el signo, el
primer miembro de (16.8.8.2); de donde se sigue la conclusión.

\begin{proposition}[16.8.9]
\label{IV.16.8.9}
Si $D:\mathscr{F}\to\mathscr{G}$ es un operador diferencial de orden
$\leq n$, y $D':\mathscr{G}\to\mathscr{H}$ un operador diferencial de orden
$\leq n'$, entonces $D'\circ D:\mathscr{F}\to\mathscr{H}$ es un operador
diferencial de orden $\leq n+n'$.
\end{proposition}

Por hipótesis, se puede escribir $D=u\circ d_{X/S,\mathscr{F}}^n$ y
$D'=v\circ d_{X/S,\mathscr{G}}^{n'}$, donde
$u:\mathscr{P}_{X/S}^n\otimes_{\mathcal{O}_X}\mathscr{F}\to\mathscr{G}$ y
$v:\mathscr{P}_{X/S}^{n'}\otimes_{\mathcal{O}_X}\mathscr{G}\to\mathscr{H}$
son $\mathcal{O}_X$-homomorfismos. Todo se reduce a demostrar que el
homomorfismo compuesto de haces de grupos aditivos
\[
  \xymatrix{
    \mathscr{F}\ar[r]^-{d_{X/S,\mathscr{F}}^n}
      &\mathscr{P}_{X/S}^n\otimes_{\mathcal{O}_X}\mathscr{F}
        \ar[r]^-u
      &\mathscr{G}\ar[r]^-{d_{X/S,\mathscr{G}}^{n'}}
      &\mathscr{P}_{X/S}^{n'}\otimes_{\mathcal{O}_X}\mathscr{G}
  }
\]
se factoriza como
\[
  \xymatrix{
    \mathscr{F}\ar[r]^-{d_{X/S,\mathscr{F}}^{n+n'}}
      &\mathscr{P}_{X/S}^{n+n'}\otimes_{\mathcal{O}_X}\mathscr{F}
        \ar[r]^-w
      &\mathscr{P}_{X/S}^{n'}\otimes_{\mathcal{O}_X}\mathscr{G}
  }
\]
donde $w$ es un $\mathcal{O}_X$-homomorfismo. Bastará demostrar el siguiente

\begin{lemma}[16.8.9.1]
\label{IV.16.8.9.1}
Existe un único $\mathcal{O}_X$-homomorfismo
\[
  \delta:\mathscr{P}_{X/S}^{n+n'}
  \longrightarrow
  \mathscr{P}_{X/S}^{n'}(\mathscr{P}_{X/S}^n)
  =\mathscr{P}_{X/S}^{n'}\otimes_{\mathcal{O}_X}\mathscr{P}_{X/S}^n
  \tag{16.8.9.2}
  \label{IV.16.8.9.2}
\]
\oldpage[IV-4]{44}
que hace conmutativo el diagrama
\[
  \xymatrix{
    \mathcal{O}_X\ar[r]^-{d_{X/S}^{n+n'}}\ar[d]_-{d_{X/S}^n}
      &\mathscr{P}_{X/S}^{n+n'}\ar[d]^-\delta\\
    \mathscr{P}_{X/S}^n
      \ar[r]_-{d_{X/S,\mathscr{P}_{X/S}^n}^{n'}}
      &\mathscr{P}_{X/S}^{n'}(\mathscr{P}_{X/S}^n)
  }
  \tag{16.8.9.3}
  \label{IV.16.8.9.3}
\]
\end{lemma}

Se tendrá entonces, en efecto, un diagrama conmutativo deducido de (16.8.9.3)
por tensorización con $\mathscr{F}$
\[
  \xymatrix{
    \mathscr{F}\ar[r]^-{d_{X/S,\mathscr{F}}^{n+n'}}
      \ar[d]_-{d_{X/S,\mathscr{F}}^n}
      &\mathscr{P}_{X/S}^{n+n'}(\mathscr{F})
        \ar[d]^-{\delta\otimes1}\\
    \mathscr{P}_{X/S}^n(\mathscr{F})
      \ar[r]_-{d_{X/S,\mathscr{P}_{X/S}^n(\mathscr{F})}^{n'}}
      &\mathscr{P}_{X/S}^{n'}(\mathscr{P}_{X/S}^n(\mathscr{F}))
  }
\]
y, por otra parte, se verifica inmediatamente a partir de la definición
(16.7.5) que el diagrama
\[
  \xymatrix{
    \mathscr{P}_{X/S}^n(\mathscr{F})\ar[r]^-u
      \ar[d]_-{d_{X/S,\mathscr{P}_{X/S}^n(\mathscr{F})}^{n'}}
      &\mathscr{G}\ar[d]^-{d_{X/S,\mathscr{G}}^{n'}}\\
    \mathscr{P}_{X/S}^{n'}(\mathscr{P}_{X/S}^n(\mathscr{F}))
      \ar[r]_-{1\otimes u}
      &\mathscr{P}_{X/S}^{n'}(\mathscr{G})
  }
\]
es conmutativo. Se resuelve, por tanto, la cuestión tomando como $w$ el
$\mathcal{O}_X$-homomorfismo compuesto
\[
  \xymatrix{
    \mathscr{P}_{X/S}^{n+n'}(\mathscr{F})
      \ar[r]^-{\delta\otimes1}
      &\mathscr{P}_{X/S}^{n'}(\mathscr{P}_{X/S}^n(\mathscr{F}))
        \ar[r]^-{1\otimes u}
      &\mathscr{P}_{X/S}^{n'}(\mathscr{G}).
  }
\]

Queda por demostrar el lema (16.8.9.1). Teniendo en cuenta (16.7.6), que
demuestra la unicidad de $\delta$, se reduce al caso en que
$S=\operatorname{Spec}(A)$ y $X=\operatorname{Spec}(B)$ son afines; poniendo
$\mathfrak{I}=\mathfrak{I}_{B/A}$, se trata de definir un homomorfismo
canónico de $B$-módulos
\[
  \varphi:(B\otimes_A B)/\mathfrak{I}^{n+n'+1}
  \longrightarrow
  ((B\otimes_A B)/\mathfrak{I}^{n'+1})
  \otimes_B((B\otimes_A B)/\mathfrak{I}^{n+1})
\]
donde las estructuras de $B$-módulo de ambos miembros provienen del
primer factor $B$; recordemos que, en el producto tensorial del segundo
miembro, $(B\otimes_A B)/\mathfrak{I}^{n'+1}$ debe considerarse
\oldpage[IV-4]{45}
como $B$-módulo a la derecha mediante su segundo factor $B$, y
$(B\otimes_A B)/\mathfrak{I}^{n+1}$ como $B$-módulo a la izquierda
mediante su primer factor $B$ (16.7.2). Esto equivale a definir un
homomorfismo de $B$-módulos
\[
  \varphi_0:B\otimes_A B\longrightarrow
  ((B\otimes_A B)/\mathfrak{I}^{n'+1})
  \otimes_B((B\otimes_A B)/\mathfrak{I}^{n+1})
\]
y demostrar que se anula sobre $\mathfrak{I}^{n+n'+1}$. Ahora bien, se
define inmediatamente tal homomorfismo mediante la condición
\[
  \varphi_0(b\otimes b')
  =\pi_{n'}(b\otimes1)\otimes\pi_n(1\otimes b')
  \qquad\text{para $b,b'$ en $B$}
\]
con las notaciones de (16.3.7). Además, es inmediato que $\varphi_0$ es un
homomorfismo de \emph{anillos}. Ahora bien, se puede escribir
\[
  \begin{aligned}
  \varphi_0(b\otimes1-1\otimes b)
  ={}&\pi_{n'}(b\otimes1-1\otimes b)\otimes\pi_n(1\otimes1)\\
  &+\pi_{n'}(1\otimes b)\otimes\pi_n(1\otimes1)
  -\pi_{n'}(1\otimes1)\otimes\pi_n(1\otimes b)
  \end{aligned}
\]
y se tiene
\[
  \begin{aligned}
  \pi_{n'}(1\otimes b)\otimes\pi_n(1\otimes1)
  &=\pi_{n'}(1\otimes1)b\otimes\pi_n(1\otimes1)\\
  &=\pi_{n'}(1\otimes1)\otimes b\pi_n(1\otimes1)\\
  &=\pi_{n'}(1\otimes1)\otimes\pi_n(b\otimes1)
  \end{aligned}
\]
de donde, finalmente,
\[
  \begin{aligned}
  \varphi_0(b\otimes1-1\otimes b)
  ={}&\pi_{n'}(b\otimes1-1\otimes b)\otimes\pi_n(1\otimes1)\\
  &+\pi_{n'}(1\otimes1)\otimes\pi_n(b\otimes1-1\otimes b).
  \end{aligned}
  \tag{16.8.9.4}
  \label{IV.16.8.9.4}
\]
Un producto de $n+n'+1$ términos de la forma (16.8.9.4) es, por tanto,
necesariamente nulo, pues así sucede con un producto de $n+1$ términos de la
forma $\pi_n(b\otimes1-1\otimes b)$ y con un producto de $n'+1$ términos de
la forma $\pi_{n'}(b\otimes1-1\otimes b)$. La conclusión resulta, por tanto,
de (0, 20.4.4).

\begin{corollary}[16.8.10]
\label{IV.16.8.10}
El haz $\mathscr{D}\!iff_{X/S}(\mathcal{O}_X,\mathcal{O}_X)$ (denotado también
por $\mathscr{D}\!iff_{X/S}$) está provisto canónicamente de una estructura de
haz de anillos, y los $\mathscr{D}\!iff_{X/S}^n$ forman una filtración
creciente compatible con esta estructura.
\end{corollary}

En particular, $\mathscr{D}\!iff_{X/S}^0$ es un haz de subanillos de
$\mathscr{D}\!iff_{X/S}$ que se identifica canónicamente con $\mathcal{O}_X$
(16.8.1). Las fórmulas (16.8.5.1) y (16.8.5.2) muestran que la estructura de
$\mathcal{O}_X$-bimódulo de $\mathscr{D}\!iff_{X/S}$ proviene de la
multiplicación por la izquierda y por la derecha mediante las secciones de
$\mathcal{O}_X$, considerado como haz de subanillos de
$\mathscr{D}\!iff_{X/S}$.

\begin{remarks}[16.8.11]
\label{IV.16.8.11}
\begin{enumerate}
  \item[(i)] Supongamos que
  $\mathscr{F}=\bigoplus_{\lambda\in L}\mathscr{F}_\lambda$; entonces es
  claro por (16.7.2.1) que
  $\mathscr{P}_{X/S}^n(\mathscr{F})=
  \bigoplus_{\lambda\in L}\mathscr{P}_{X/S}^n(\mathscr{F}_\lambda)$; como el
  funtor $\mathscr{F}\mapsto\Gamma(U,\mathscr{F})$ conmuta con la formación
  de sumas directas cualesquiera, $d_{X/S,\mathscr{F}}^n$ es el homomorfismo
  cuya restricción a cada $\mathscr{F}_\lambda$ es
  \[
    d_{X/S,\mathscr{F}_\lambda}^n:\mathscr{F}_\lambda
    \longrightarrow\mathscr{P}_{X/S}^n(\mathscr{F}_\lambda);
  \]
  se concluye inmediatamente que se tiene
  \[
    \operatorname{Diff}_{X/S}^n(\mathscr{F},\mathscr{G})
    =\prod_{\lambda\in L}
    \operatorname{Diff}_{X/S}^n(\mathscr{F}_\lambda,\mathscr{G}),
  \]
  y, por consiguiente, también (0\textsubscript{I}, 3.2.6)
  \[
    \mathscr{D}\!iff_{X/S}^n(\mathscr{F},\mathscr{G})
    =\prod_{\lambda\in L}
    \mathscr{D}\!iff_{X/S}^n(\mathscr{F}_\lambda,\mathscr{G}).
  \]
  Por otra parte, si $\mathscr{G}=\prod_{\mu\in M}\mathscr{G}_\mu$
  (0\textsubscript{I}, 3.2.6), se tiene
  \[
    \operatorname{Hom}_{\mathcal{O}_X}
    (\mathscr{P}_{X/S}^n(\mathscr{F}),\mathscr{G})
    =\prod_{\mu\in M}\operatorname{Hom}_{\mathcal{O}_X}
    (\mathscr{P}_{X/S}^n(\mathscr{F}),\mathscr{G}_\mu),
  \]
  \oldpage[IV-4]{46}
  pues todo homomorfismo $u$ de $\mathscr{P}_{X/S}^n(\mathscr{F})$ en
  $\mathscr{G}$ corresponde biunívocamente a la familia de sus compuestos
  \[
    u_\mu:\mathscr{P}_{X/S}^n(\mathscr{F})
    \longrightarrow\mathscr{G}\longrightarrow\mathscr{G}_\mu.
  \]
  Se tiene, por tanto,
  \[
    \operatorname{Diff}_{X/S}^n(\mathscr{F},\mathscr{G})
    =\prod_{\mu\in M}
    \operatorname{Diff}_{X/S}^n(\mathscr{F},\mathscr{G}_\mu),
  \]
  y, por consiguiente, también
  \[
    \mathscr{D}\!iff_{X/S}^n(\mathscr{F},\mathscr{G})
    =\prod_{\mu\in M}
    \mathscr{D}\!iff_{X/S}^n(\mathscr{F},\mathscr{G}_\mu).
  \]
  \item[(ii)] Hasta aquí apenas se han encontrado operadores diferenciales
  $\mathscr{F}\to\mathscr{G}$ salvo cuando $\mathscr{F}$ y $\mathscr{G}$ son
  $\mathcal{O}_X$-módulos localmente libres de rango finito, en cuyo caso, en
  virtud de (i), su estructura se reduce localmente a la del haz
  $\mathscr{D}\!iff_{X/S}$; esta última se estudiará más adelante (16.11) en
  un caso particular.
\end{enumerate}
\end{remarks}
